\documentclass{beamer}
%%%%%%%%%%%%%%%%%%%%%%%%%%%%%%%%%%%%%%%%%%%%%%%%%%%%%%%%%%%%%%%%%%%%%%%%%%%%%%%%%%%%%%%%%%%%%%%%%%%%%%%%%%%%%%%%%%%%%%%%%%%%%%%%%%%%%%%%%%%%%%%%%%%%%%%%%%%%%%%%%%%  DATOS DE LA PLANTILLA %%%%%%%%%%%%%%%%%%%%%%%%%% %%%%%%%%%%%%%%%%%%%%%%%%%%%%%%%%%%%%%%%%%%%%%%%%%%%%%%%%%%%%%%%%%%%%%%%%%%%%%%%%%%%%%%%%%%%%%%%%%%%%%%%%%%%%%%%%%%%%%%%%%%%%%%%%



\useinnertheme[shadow=true]{rounded}

\usepackage{etoolbox}

\setbeamercolor{block title}{use=structure,fg=structure.fg,bg=structure.fg!20!bg}
\setbeamercolor{block body}{parent=normal text,use=block title,bg=block title.bg!50!bg}

\setbeamercolor{block title example}{use=example text,fg=example text.fg,bg=example text.fg!20!bg}
\setbeamercolor{block body example}{parent=normal text,use=block title example,bg=block title example.bg!50!bg}

\addtobeamertemplate{proof begin}{%
    \setbeamercolor{block title}{fg=black,bg=red!50!white}
    \setbeamercolor{block body}{fg=red, bg=red!30!white}
}{}

\BeforeBeginEnvironment{theorem}{
    \setbeamercolor{block title}{fg=black,bg=orange!50!white}
    \setbeamercolor{block body}{fg=orange, bg=orange!30!white}
}
\AfterEndEnvironment{theorem}{
 \setbeamercolor{block title}{use=structure,fg=structure.fg,bg=structure.fg!20!bg}
 \setbeamercolor{block body}{parent=normal text,use=block title,bg=block title.bg!50!bg, fg=black}
}

\BeforeBeginEnvironment{definition}{%
    \setbeamercolor{block title}{fg=black,bg=pink!50!white}
    \setbeamercolor{block body}{fg=pink, bg=pink!30!white}
}
\AfterEndEnvironment{definition}{
 \setbeamercolor{block title}{use=structure,fg=structure.fg,bg=structure.fg!20!bg}
 \setbeamercolor{block body}{parent=normal text,use=block title,bg=block title.bg!50!bg, fg=black}
}





%%%%%%%%%%%%%%%%%%%%%%%%%%%%%%%%%%%%%%%%%%%%%%%%%%%%%%%%%%%%%%%%%%%%%%%%%%%%%%%%%%%%%%%%%%%%%%%%%%%%%%%%%%%%%%%%%%%%%%%%%%%%%%%%%%%%%%%%%%%%%%%%%%%%%%%%%%%%%%%%%%%%%%%%%%%%%%%%%%%%%%%%%%%%%%%%%%%%%%%%%%%%%%%%%%%%%%%%%%%%%%%%%%%%%%%%%%%%%%%%%%%%%%%%%%%%%%%%%%%%%%%%%%%%%%%%%%%%%%%%%%%%%%%%%%%%%%%%%%%%%%%%%%%%%%%%%%%%%%%%%%%%%%%%%%%%%%%%%%%%%%%%%%%%%%%%%%%%%%%%%%%%%%%%%%%%%%%%%%%%%%%%%%%%%%%%%%%%%%%%%%%%%%%%%%%%%%%%%%%%%%%%%%%%%%%%%%%%%%%%%%%%%%%%
\usetheme[titlepagelogo=fcfm,% Logo for the first page
		  language=spanish,
		  coding=utf8,
		  bullet=triangle,
		  color=blue,%green, blue
		  notshowauthor=true,%Para no mostrar el nombre del autor en cada frame
		  notshowtitle=true,
         ]{TorinoTh}
         
\usepackage[beamer,customcolors]{hf-tikz}
\setbeamercovered{transparent}
\hfsetfillcolor{alerted text.fg!10}
\hfsetbordercolor{alerted text.fg}
%%%%%%%%%%%%%%%%%%%%%%%%%%%%%%%%%%%%%%%%%%%%%%%%%%%%%%%%%%%%%%%%%%%%%%%%%%%%%%%%%%%%%%%%%%%%%%%%%%%%%%%%%%%%%%%%%%%%%%%%%%%%%%%%%%%%%%%%%%%%%%%%%%%%%%%%%%%%%%%%%%%  DATOS DE LA PORTADA %%%%%%%%%%%%%%%%%%%%%%%%%% %%%%%%%%%%%%%%%%%%%%%%%%%%%%%%%%%%%%%%%%%%%%%%%%%%%%%%%%%%%%%%%%%%%%%%%%%%%%%%%%%%%%%%%%%%%%%%%%%%%%%%%%%%%%%%%%%%%%%%%%%%%%%%%%
\author{\footnotesize Roque Vidal Luciano Gerardo}

\rel{\footnotesize Dr. Juan Alberto Escamilla Reyna\\ Dr. Tomás Pérez Becerra}
\title{``Teoremas de Punto Fijo en Espacios Normados y sus Aplicaciones en Ecuaciones Integrales"}
\ateneo{FCFM, BUAP}
\setcandidatelabel{Presenta}
\setrellabel{Asesores}
\date{2 de Diciembre de 2025}
\thispagestyle{empty}%%Este comando sirve para quitar la numeración de esta página (la de la portada). 


%%%%%%%%%%%%%%%%%%%%%%%%%%%%%%%%%%%%%%%%%%%%%%%%%%%%%%%%%%%%%%%%%%%%%%%%%%%%%%%%%%%%%%%%%%%%%%%%%%%%%%%%%%%%%%%%%%%%%%%%%%%%%%%%%%%%%%%%%%%%%%%%%%%%%%%%%%%%%%%%%%%  LIBRERÍAS USADAS %%%%%%%%%%%%%%%%%%%%%%%%%% %%%%%%%%%%%%%%%%%%%%%%%%%%%%%%%%%%%%%%%%%%%%%%%%%%%%%%%%%%%%%%%%%%%%%%%%%%%%%%%%%%%%%%%%%%%%%%%%%%%%%%%%%%%%%%%%%%%%%%%%%%%%%%%%
\usepackage[utf8]{inputenc} 
\usepackage{tikz} 
\usepackage{amsmath}
\usepackage{sidecap}
\usepackage{wrapfig}
\usepackage{soul}
\usepackage{microtype}
\usepackage{array}

%\usepackage{makeidx}
\usetikzlibrary{positioning,decorations.pathreplacing,shapes}

\usetikzlibrary{shapes,arrows} 
 \usetikzlibrary{shapes,backgrounds}
%Definir colores
\definecolor{Azul}{rgb}{0.1,0.1,0.6}


%%%%%%%%%%%%%%%%%%%%%%%%%%%%%%%%%%%%%%%%%%%%%%%%%%%%%%%%%%%%%%%%%%%%%%%%%%%%%%%%%%%%%%%%%%%%%%%%%%%%%%%%%%%%%%%%%%%%%%%%%%%%%%%%%%%%%%%%%%%%%%%%%%%%%%%%%  INICIO DE DOCUMENTO    %%%%%%%%%%%%%%%%%%%%%%%%%%%%%%%%% %%%%%%%%%%%%%%%%%%%%%%%%%%%%%%%%%%%%%%%%%%%%%%%%%%%%%%%%%%%%%%%%%%%%%%%%%%%%%%%%%%%%%%%%%%%%%%%%%%%%%%%%%%%%


\begin{document}

	\maketitle
	%\setbeamercovered{dynamic}
	%\titlepageframe
	
	\begin{frame}
    \frametitle{Índice}
    \tableofcontents

    \thispagestyle{empty}
\end{frame}




%%%%%%%%%%%%%%%%%%%%%%%%%%%%%%%%%%%%%%%%%%%%%%%%%%%%%%%%%%%%%%%%%%%%%%%%%%%%%%%%%%%%%%%%%%%%%%%%%%%%%%%%%%%%%%%%%%%%%%%%%%%%%%%%%%%%%%%%%%%%%%%%%%%%%%%%%%%%%%%%%%%%%%%%%%%%%%%%%%%%%%%%%%%%%%%%%%%%%%%%%%%%%%%%%%%%%%%%%%%%%%%%%%%%%%%%%%%%%%%%%%%%%%%%%%%%%%%%%%%%%%%%%%%%%%%%%%%%%%%%%%%%%%%%%%%%%%%%%%%%%%%%%%%%%%%%%%%%%%%%%%%%%%%%%%%%%%


%\section{\textcolor{blue}{Actividades realizadas durante el semestre}}
\begin{frame}{}
    \begin{block}{}
	\centering
	\LARGE  \textcolor{Azul}{Actividades realizadas durante el semestre} 
	
\end{block}


\begin{itemize}
	\item \textbf{Materias que se están cursando:}
	\begin{itemize}
	   \item[$ \ast $] SEM. DE TESIS E INVESTIGACIÓN VI 
	\end{itemize}
	\item \textbf{Participaciones}
	\begin{itemize}
		\item[$ \ast $] 	Seminario del Cuerpo Académico de Análisis Matemático,\\ 
         1 de Marzo de 2023
	\end{itemize}
	
\end{itemize}


\end{frame}

%%%%%%%%%%%%%%%%%%%%%%%%%%%%%%%%%%%%%%%%%%%%%%%%%%%%%%%%%%%%%%%%%%%%%%%%%%%%%%%%%%%%%%%%%%%%%%%%%%%%%%%%%%%%%%%%%%%%%%%%%%%%%%%%%%%%%%%%%%%%%%%%%%%%%%%%%%%%%%%%%%%%%%%%%%%%%%%%%%%%%%%%%%%%%%%%%%%%%%%%%%%%%%%%%%%%%%%%%%%%%%%%%%%%%%%%%%%%%%%%%%%%%%%%%%%%%%%%%%%%%%%%%%%%%%%%%%%%%%%%%%%%%%%%%%%%%%%%%%%%%%%%%%%%%%%%%%%%%%%%%%%%%%%%%%%%%%





%\section{\textcolor{blue}{Actividades realizadas durante el semestre}}
%\begin{frame}{}
 %   \begin{block}{}
%	\centering
%	\LARGE  \textcolor{Azul}{Actividades realizadas durante el semestre} 
	
%\end{block}


%\begin{itemize}
%	\item \textbf{Congresos:}
%	\begin{itemize}
%		\item[$ \ast $] 	57 CONGRESO NACIONAL DE LA SOCIEDAD MATEMÁTICA MEXICANA,\\ 
 %        21-25 de Octubre de 2024
%	\end{itemize}
%	\item \textbf{Materias impartidas:}
%	\begin{itemize}
%	   \item[$ \ast $]  Matemáticas Básicas para Física, FCFM, BUAP.\\
        %7 de Agosto de 2023 al 7 de Diciembre de 2023
%	\end{itemize}
%\end{itemize}


%\end{frame}



%%%%%%%%%%%%%%%%%%%%%%%%%%%%%%%%%%%%%%%%%%%%%%%%%%%%%%%%%%%%%%%%%%%%%%%%%%%%%%%%%%%%%%%%%%%%%%%%%%%%%%%%%%%%%%%%%%%%%%%%%%%%%%%%%%%%%%%%%%%%%%%%%%%%%%%%%%%%%%%%%%%%%%%%%%%%%%%%%%%%%%%%%%%%%%%%%%%%%%%%%%%%%%%%%%%%%%%%%%%%%%%%%%%

\section{\textcolor{blue}{Preliminares: ENO}}

\begin{frame}{Preliminares: ¿Qué es un Espacio Normado Ordenado?}
	\begin{columns}
		\column{.45\textwidth}
		$-$\textbf{Definición: \cite{cone2007} }
		\begin{block}{}
			Sea $ (X, || \cdot || ) $ un espacio normado y  $ X_{+} \subset  X$ cerrado no vacío.\\ Se dice que $ X_{+} $ es un \textbf{cono}\\ si $ x, \ y \in X_{+} $, $\lambda \geq 0$ cumplen:
			\begin{itemize}
				%\item $P$ es cerrado en $X$.
				\item  $x + y \in X_{+}$.
				\item  $\lambda x \in X_{+}$.
				\item  $-x, \ x \in X_{+}$ implica $x = \bar{0} $.
				%\item $P$ es cerrado.
			\end{itemize}
			
		\end{block}
		A  $(X, || \cdot | | , X_{+}  )$ se le llama \textbf{espacio normado ordenado}.\\
		
		\column{.45\textwidth}
	
			\begin{tikzpicture}
			%Tetraedro azul
			\draw[color=blue] (0,-1) -- (-1,1) -- (1,1) -- (0,-1); 
			\draw[color=blue] (-1,1) -- (0,1.5) -- (1,1);
			\draw[ dashed, color=blue]  (0,-1) -- (0,1.5);
			%Tetraedro Rojo
			\draw[color=red] (0,-1) -- (-1,-3) -- (1,-3) -- (0,-1); 
			\draw[ dashed, color=red] (-1,-3) -- (0,-2.5) -- (1,-3);
			\draw[ dashed, color=red]  (0,-1) -- (0,-2.5);
		\end{tikzpicture}
		\textcolor{blue}{$ X_{+} + \lambda X_{+} \subset X_{+} $.}\\ \textcolor{blue}{$X_{+} $} $\cap$  \textcolor{red}{$(- x_{+} )$} $= \{ \bar{0} \}$.
	\end{columns} 
\end{frame}

%%%%%%%%%%%%%%%%%%%%%%%%%%%%%%%%%%%%%%%%%%%%%%%%%%%%%%%%%%%%%%%%%%%%%%%%%%%%%%%%%%%%%%%%%%%%%%%%%%%%%%%%%%%%%%%%%%%%%%%%%%%%%%%%%%%%%%%%%%%%%%%%%%%%%%%%%%%%%%%%%%%%%%%%%%%%%%%%%%%%%%%%%%%%%%%%%%%%%%%%%%%%%%%%%%%%%%%%%%%%%%%%%%%%%%%%%%%%%%%%%%%%%%%%%%%%%%%%%%%%%%%%%%%%%%%%%%%%%%%%%%%%%%%%%%%%%%%%%%%%%%%%%%%%%%%%%%%%%%%%%%%%%%%%%%%%%%%%%%%%%%%%%%%%%%%%%%%%%%%%%%%%%%%%%%%%


\begin{frame}{}
	\begin{block}{Teorema}
			Sean $(X, || \cdot || , X_{+} )$ un ENO, $ a,b,c \in X$ y $ \lambda \geq 0$. Entonces la relación  $``\leq "$ en $X$   definida como $a \leq b$ si $b-a \in X_{+}$  cumple que
		\begin{enumerate}[(1)]
			\item $a \leq a$, 
			\item Si $a \leq b$ y $b \leq a$, entonces $ a=b $, 
			\item Si $a \leq b$ y $b \leq c$, entonces $a \leq c$, 
			\item Si $a \leq b$, entonces $a+c \leq b+c$, 
			\item $ a \leq b$, entonces $\lambda a \leq \lambda b $.
		\end{enumerate}
	\end{block}
Nota:	$ a < b $, si $ a \leq b  $ y  $ a \neq b . $ 
\end{frame}
	%%%%%%%%%%%%%%%%%%%%%%%%%%%%%%%%%%%%%%%%%%%%%%%%%%%%%%%%%%%%%%%%%%%%%%%%%%%%%%%%%%%%%%%%%%%%%%%%%%%%%%%%%%%%%%%%%%%%%%%%%%%%%%%%%%%%%%%%%%%%%%%%%%%%%%%%%%%%%%%%%%%%%%%%%%%%%%%%%%%%%%%%%%%%%%%%%%%%%%%%%%%%%%%%%%%%%%%%%%%%%%%%%%%%%%%%%%%%%%%%%%%%%%%%%%%%%%%%%%%%%%%%%%%%%%%%%%%%%%%%%%%%%%%%%%%%%%%%%%%%
	\begin{frame}{}
		$-$\textbf{Definición:} 
		\begin{block}{}
		    	Sean $(X,|| \cdot || , X_{+} ) $ un ENO y $ a, b \in X_{+} $ con $ a \leq b  $. 
		\begin{itemize}
			\item $[a , b ] = \left\lbrace  x \in X : a \leq x \leq b \right\rbrace  $ se llama  \textbf{intervalo ordenado cerrado} de $(X,\leq)$.
	 	\item $[a, \infty ) = \left\lbrace  x \in X : a \leq x  \right\rbrace  $  se llama  \textbf{intervalo ordenado infinito  semicerrado por la izquierda } de $(X,\leq)$.\\
    Análogamente se define $( - \infty, a ] $ el \textbf{intervalo ordenado infinito  semicerrado por la derecha}.
		\end{itemize}
		\end{block}
		
	\end{frame}


	
%%%%%%%%%%%%%%%%%%%%%%%%%%%%%%%%%%%%%%%%%%%%%%%%%%%%%%%%%%%%%%%%%%%%%%%%%%%%%%%%%%%%%%%%%%%%%%%%%%%%%%%%%%%%%%%%%%%%%%%%%%%%%%%%%%%%%%%%%%%%%%%%%%%%%%%%%%%%%%%%%%%%%%%%%%%%%%%%%%%%%%%%%%%%%%%%%%%%%%%%%%%%%%%%%%%%%%%%%%%%%%%%%%%%%%%%

%% ─────────────────────────────────────────────────────────────────────────────
%%  Clasificación de conos y ejemplos de ENO
%% ─────────────────────────────────────────────────────────────────────────────
\section{Clasificación de conos}

\begin{frame}{Tipos de conos en un ENO}
\framesubtitle{Propiedades geométricas que determinan el comportamiento de BV}

\begin{block}{}
\centering
\large Algunos tipos de conos
\end{block}

\begin{itemize}
    \item \textbf{Conos normales}
    \begin{itemize}
        \item[$\ast$] \textbf{Normal:} existe $M>0$ tal que
            $0 \le x \le y \Rightarrow \|x\| \le M\|y\|$.
        \item[$\ast$] \textbf{Extendible:} existe $x^{*}\in \Xp^{*}$ y $\alpha>0$
            tal que $x^{*}(x)\ge\alpha\|x\|$ para todo $x\in\Xp$.
    \end{itemize}
    \item \textbf{Conos no normales}
    \begin{itemize}
        \item[$\ast$] \textbf{Generador:} $X = \Xp - \Xp$, es decir, todo
            $x\in X$ se escribe como $x=\alpha-\beta$ con $\alpha,\beta\in\Xp$.
        \item[$\ast$] \textbf{Sólido:} $\operatorname{int}(\Xp)\ne\emptyset$.
    \end{itemize}
\end{itemize}

% IDEA ORAL: Recordar que estas propiedades no son técnicas por sí mismas,
% sino que cada una tiene una consecuencia concreta sobre las funciones BV.

\end{frame}

%──────────────────────────────────────────────────────────────────────────────

\begin{frame}{Relaciones entre tipos de conos}
\framesubtitle{¿Qué implica qué?}

\begin{columns}[c]
\column{.60\textwidth}
\begin{tikzpicture}
    \begin{scope}[shift={(3cm,-3cm)}, fill opacity=0.5]
        \def\Gcirc{(0.5,0) circle (2cm)}
        \def\Ncirc{(0:3.7cm) circle (2cm)}
        \fill[green!60!white] \Gcirc;
        \fill[red!60!white]   \Ncirc;
        \draw (-2,-2.5) rectangle (7,2)
              node[text=black,above,fill opacity=1,text opacity=1]
              {$(X,\|\cdot\|,\Xp)$ ENO};
    \end{scope}
\end{tikzpicture}

\column{.38\textwidth}
\begin{itemize}
    \item \textcolor{green!70!black}{Verde:} cono generador
    \item \textcolor{red!70!black}{Rojo:} cono normal
    \item \textcolor{brown}{Intersección:} generador y normal
\end{itemize}

\medskip
\small
\textbf{Sólido} $\Rightarrow$ generador,
pero no al revés.

\end{columns}

\end{frame}

%──────────────────────────────────────────────────────────────────────────────

\begin{frame}{Ejemplo 1: ENO con cono normal}
\framesubtitle{$C^1_{\mathbb{R}}[a,b]$ con la norma del supremo}

\begin{block}{}
Sea $C_{\mathbb{R}}^{1}[a,b]$ con
$\|f\|_{\infty} = \sup_{x\in[a,b]}|f(x)|$
y el cono $\Xp = \{f\in C_{\mathbb{R}}^{1}[a,b] : f(x)\ge 0\}$.

$(C_{\mathbb{R}}^{1}[a,b],\|\cdot\|_{\infty},\Xp)$ es un ENO y $\Xp$ es \textbf{normal}.
\end{block}

\begin{columns}
\column{.58\textwidth}
\begin{tikzpicture}[scale=0.85]
    \draw[->] (-3.5,0) -- (3.9,0) node[below] {$x$};
    \draw[->] (0,-0.5) -- (0,3.9) node[left] {$y$};
    \draw[smooth,<->,domain=-3.16:3.17,color=green!70!black,very thick]
        plot (\x,{sin(\x r)*(\x+1)+1});
    \node at (2,3.9) {\small\textcolor{green!70!black}{$f\ge 0$}};
    \draw[fill] (-3.16,0) circle (0.07) node[below]{$a$};
    \draw[fill] (3.3,0) circle (0.07) node[below]{$b$};
\end{tikzpicture}

\column{.40\textwidth}
\textbf{Consecuencia:} \\[4pt]
$0\le f\le g \Rightarrow \|f\|_\infty\le\|g\|_\infty$.

\medskip
Las sucesiones crecientes acotadas convergen.
\end{columns}

\end{frame}

%──────────────────────────────────────────────────────────────────────────────

\begin{frame}{Ejemplo 2: ENO con cono \textbf{no} normal}
\framesubtitle{$C^1_{\mathbb{R}}[0,1]$ con la norma $\|\cdot\|_\infty+\|\cdot'\|_\infty$}

\begin{block}{}
Sea $C_{\mathbb{R}}^{1}[0,1]$ con
$\|f\| = \|f\|_\infty + \|f'\|_\infty$
y $\Xp = \{f\in C_{\mathbb{R}}^{1}[0,1]: f(x)\ge 0\}$.

$\Xp$ \textbf{no es normal}.
\end{block}

\begin{columns}
\column{.55\textwidth}
\begin{tikzpicture}[scale=0.8]
    \draw[->] (-3.5,0) -- (3.9,0) node[below]{$x$};
    \draw[->] (0,-0.5) -- (0,3.9) node[left]{$y$};
    \draw[smooth,<->,domain=-3.16:3.17,color=green!70!black,very thick]
        plot (\x,{sin(\x r)+1});
    \node at (-0.5,0.5){\small\textcolor{green!70!black}{$f_n$}};
    \draw[->,color=MyTeal] (-3.16,2) -- (3.3,2) node[below]{$\bar{2}$};
    \draw[fill] (-3.16,0) circle (0.07) node[below]{$0$};
    \draw[fill] (3.3,0) circle (0.07) node[below]{$1$};
\end{tikzpicture}

\column{.43\textwidth}
$f_n(t)=\sin(2\pi n t)+1$.

\medskip
$0\le f_n\le\bar{2}$, pero
\[
\|f_n\| = 1 + 2\pi n \to\infty.
\]

La norma no está controlada por el orden.
\end{columns}

\end{frame}

   
	

\begin{frame}{Espacios de funciones vectoriales sobre ENO's }
\begin{block}{}
    Sean $(X,|| \cdot || , X_{+} ) $ un ENO. Una función  \textcolor{green}{$ f: [a,b] \longrightarrow X $} se dice que es una función vectorial.
\end{block}
    \begin{tikzpicture}
	    % Dibujo el Dominio					+	+	+
      \draw[color=blue] (0,0) -- (3,0);
      % Codominio  (Rectángulo ):
      \draw[fill=gray] (5,-1) rectangle (9,2) node[fill=red!20,draw,double,rounded corners] {$ X $};
     
    %  Partición 
     \draw[fill] (0,0) circle (0.07) node[below]{$ a $}; 
     % \draw[fill] (1,0) circle (0.07) node[below]{$t_{i-1}  $};
        %\draw[fill] (1.5,0) circle (0.07) node[below]{$t_{i}  $}; 
          \draw[fill] (3,0) circle (0.07) node[below]{$ b  $}; 
        %f(x_{i} ) , f(x_{i-1} )
      
     
        % punto inicial, grado inicial, final, radio
\draw[->,color=green] (2,0) to[bend left] (7,1.4);

\node (1) at (3.5,2.3){ \textcolor{green}{$f$}};
	\end{tikzpicture}


\end{frame}




%%%%%%%%%%%%%%%%%%%%%%%%%%%%%%%%%%%%%%%%%%%%%%%%%%%%%%%%%%%%%%%%%%%%%%%%%%%%%%%%%%%%%%%%%%%%%%%%%%%%%%%%%%%%%%%%%%%%%%%%%%%%%%%%%%%%%%%%%%%%%%%%%%%%%%%%%%%%%%%%%%%%%%%%%%%%%%%%%%%%%%%%%%%%%%%%%%%%%%%%%%%%%%%%%%%%%%%%%%%%%%%%%%%%%%%%%%%%%%%%%%%%%%%%%%%%%%%%%%%%%%%%%%%%%%%%%%%%%%%%%%%%%%%%%%%%%%%%%%%%



\section{V. A. débil y fuerte para funciones crecientes}
\begin{frame}{}
    \begin{center}
        {\Huge  Algunos resultados sobre ``Variación Acotada" para funciones crecientes }
    \end{center}
\end{frame}




%%%%%%%%%%%%%%%%%%%%%%%%%%%%%%%%%%%%%%%%%%%%%%%%%%%%%%%%%%%%%%%%%%%%%%%%%%%%%%%%%%%%%%%%%%%%%%%%%%%%%%%%%%%%%%%%%%%%%%%%%%%%%%%%%%%%%%%%%%%%%%%%%%%%%%%%%%%%%%%%%%%%%%%%%%%%%%%%%%%%%%%%%%%%%%%%%%%%%%%%%%%%%%%%%%%%%%%%%%%%%%%%%%%%%%%%%%%%%%%%%%%%%%%%%%%%%%%%%%%%%%%%%%%%%%%%%%%%%%%%%%%%%%%%%%%%%%%%%%%%


%\setbeamercovered{invisible}
\begin{frame}[t,fragile]{}
\textbf{Función creciente}
           \begin{block}{ }
            Sean $ (X, || \cdot || ,  P )  $ un espacio normado ordenado y $f : [a,b] \longrightarrow X  $ una función vectorial. $f$ es creciente si $ x < y $ implica $ f (x) \leq f(y)  $. 
        \end{block} 


\end{frame}

%%%%%%%%%%%%%%%%%%%%%%%%%%%%%%%%%%%%%%%%%%%%%%%%%%%%%%%%%%%%%%%%%%%%%%%%%%%%%%%%%%%%%%%%%%%%%%%%%%%%%%%%%%%%%%%%%%%%%%%%%%%%%%%%%%%%%%%%%%%%%%%%%%%%%%%%%%%%%%%%%%%%%%%%%%%%%%%%%%%%%%%%%%%%%%%%%%%%%%%%%%%%%%%%%%%%%%%%%%%%%%%%%%%%%%%%%%%%%%%%%%%%%%%%%%%%%%%%%%%%%%%%%%%%%%%%%%%%%%%%%%%%%%%%%%%%%%%%%%%%


\begin{frame}[t,fragile]{}

\textbf{Definiciones:}
\begin{itemize}	        
	\item \textbf{Variación acotada fuerte}  
 \begin{block}{}
     Se dice que $f$ \textbf{es de variación acotada sobre } $[a,b]$ si  $ Var_{a}^{b} ( f , X ) : = \sup \{ \sum_{i=1}^{n} || f(t_{i} ) - f(t_{i-1} ) || : P \in \wp  [a,b] \} < \infty $.
 
 
 \end{block}
       \begin{tikzpicture}
	    % Dibujo el Dominio					+	+	+
      \draw[color=blue] (0,0) -- (3,0);
      % Codominio  (Rectángulo ):
      \draw[fill=gray] (5,-1) rectangle (11,2) node[fill=red!20,draw,double,rounded corners] {$ X $};
      %Función 
      \draw[->,color=green] (1.5,0.2) to[bend left] (6,0.8);
      \draw[->,color=green] (1,0.2) to[bend left] (7,1.4);
    %  Partición 
     \draw[fill] (0,0) circle (0.07) node[below]{$ a $}; 
      \draw[fill] (1,0) circle (0.07) node[below]{$t_{i-1}  $};
        \draw[fill] (1.5,0) circle (0.07) node[below]{$t_{i}  $}; 
          \draw[fill] (3,0) circle (0.07) node[below]{$ b  $}; 
        %f(x_{i} ) , f(x_{i-1} )
        \shadedraw (6.1,0.8) circle (0.05cm) node[below]{$f(t_{i-1})$};
        \shadedraw (7.1,1.4) circle (0.05cm) node[below]{$f(t_{i})$};
        %Unión de f(t_{i} ) a  f(t_{i-1} )
        \draw[dashed, color=red] (6.1,0.8) -- (7.1,1.4);
	\end{tikzpicture}
    \item[$ - $]   \textbf{Variación Acotada Débil}
    \item[$ - $] \textbf{Variación en orden}
\end{itemize}

\end{frame}

%%%%%%%%%%%%%%%%%%%%%%%%%%%%%%%%%%%%%%%%%%%%%%%%%%%%%%%%%%%%%%%%%%%%%%%%%%%%%%%%%%%%%%%%%%%%%%%%%%%%%%%%%%%%%%%%%%%%%%%%%%%%%%%%%%%%%%%%%%%%%%%%%%%%%%%%%%%%%%%%%%%%%%%%%%%%%%%%%%%%%%%%%%%%%%%%%%%%%%%%%%%%%%%%%%%%%%%%%%%%%%%%%%%%%%%%%%%%%%%%%%%%%%%%%%%%%%%%%%%%%%%%%%%%%%%%%%%%%%%%%%%%%%%%%%%%%%%%%%%%

\begin{frame}[t,fragile]{}

\textbf{Definiciones:}
\begin{itemize}      
	\item[$ - $] \textbf{Variación acotada fuerte} 
  \item \textbf{Variación Acotada Débil}
     \begin{columns}
	\column{.70\textwidth}
    \begin{tikzpicture}
%Ejes x and y
    \draw[->] (-3.5,0) -- (3.9,0) node[below] {$x$};
    \draw[->] (0,-1) -- (0,3.9) node[left] {$y$};
    %Grafica verde de la función f
    \draw[smooth,<->,domain = -3.16:3.17, color=green,very thick]plot (\x,{sin(\x r)*(\x + 1 )+1 });\node at (2,3.9){\small \textcolor{green}{$f$}};
    %Partición 
    \draw[fill] (-3.16,0) circle (0.07) node[below]{$t_{0} = a  $};
    \draw[fill] (-1.8,0) circle (0.07) node[below]{$t_{1}  $};
    \draw[fill] (-0.5,0) circle (0.07) node[below]{$t_{2}   $};
    \draw[fill] (1.39,0) circle (0.07) node[below]{$t_{3}   $};
    \draw[fill] (3.3,0) circle (0.07) node[below]{$t_{4} = b  $};
  
      %Partición evaluada en f
    \draw[fill] (-3.16,1) circle (0.07) node[below]{$(\phi \circ f)(t_{0} )   $};
    \draw[fill] (-1.8,1.8) circle (0.07) node[below]{$( \phi \circ f )(t_{1} )  $};
    \draw[fill] (-0.5,0.78) circle (0.07) node[below]{$( \phi \circ f )(t_{2} )   $};
    \draw[fill] (1.39,3.38) circle (0.07) node[below]{$ ( \phi \circ f ) (t_{3} )   $};
    \draw[fill] (3.2,1) circle (0.07) node[below]{$ ( \phi \circ f ) (t_{4} )  $};
  %lineas que unen f(xi) con f(xi-1)
\draw[dashed, color= purple] (-3.16,1) -- (-1.8,1.8);
\draw[dashed, color= purple] (-1.8,1.8) -- (-0.5,0.78);
\draw[dashed, color= purple] (-0.5,0.78) -- (1.39,3.38);
\draw[dashed, color= purple] (1.39,3.38) -- (3.17,1); 
\end{tikzpicture}
	\column{.3\textwidth}

	 \begin{block}{}

      Sea $f : [a,b] \longrightarrow X  $\\ una función vectorial.\\ $f$ es de variación\\ acotada débil\\ si para todo $ \phi \in (X)^{*} $ se cumple que \textcolor{green}{$ \phi \circ f: [a,b] \longrightarrow \mathbb{R} $}\\ es de\\ \textcolor{purple}{variación acotada}.
	 \end{block}
\begin{tikzpicture}
    % Eje X
    \draw[->] (-1,0) -- (3,0) node[right] {$x$};
    % Eje Y
    \draw[->] (0,-1) -- (0,2) node[above] {$y$};

    % Función característica de [0,1]
    \draw[thick, blue] (0,0) -- (0,1) -- (1,1) -- (1,0);
    \node at (0.5, -0.5) {$\chi_{[0,1]}(x)$};

    % Función característica de [0,2]
    \draw[thick, red] (0,0) -- (0,1) -- (2,1) -- (2,0);
    \node at (1.5, -0.5) {$\chi_{[0,2]}(x)$};

    % Puntos en el eje X
    \foreach \x in {0,1,2}
        \draw (\x,1pt) -- (\x,-3pt) node[anchor=north] {$\x$};
\end{tikzpicture}
\end{columns} 
 
	
\end{itemize}

\end{frame}


\begin{frame}{}
    \begin{itemize}
        \item[$ - $] \textbf{Variación acotada fuerte} 
  \item[$ - $] \textbf{Variación Acotada Débil}
  \item \textbf{Variación Acotada en Orden}
  \begin{block}{}
      Sean $ (X,|| \cdot || ) $ un espacio espacio de Riesz. Una función  $f : [a,b] \longrightarrow X $ es de \textbf{variación acotada respecto al orden}, si existe  $M \in X $ talque,
     $$ V_{f}^{o} : = \sup \{ \sum_{i=1}^{n} | f(t_{i} ) - f(t_{i-1} ) | : P \in \wp  [a,b] \} \leq  M .$$
  \end{block}
    \end{itemize}
\end{frame}



\begin{frame}{}
    \begin{block}{Teorema}
        Sea $X$ un retículo de Banach. Entonces 
        \begin{enumerate}
            \item Toda función creciente es orden acotada.
            \item Si es orden acotada, es débil acotada. 
        \end{enumerate}
    \end{block}
\end{frame}

\begin{frame}{}
    El siguiente ejemplo es de una función creciente que  es de variación acotada débil, pero no es de variación acotada. 


 \begin{block}{Ejemplo}
   Sean  $F : [a,b] \longrightarrow L_{\infty [a,b]}  $ dada por $F(t) = \chi_{[0,t]} $ y $ ( L_{\infty [a,b]} , || \cdot ||_{L_{\infty [a,b]} } , P ) $ un ENO, donde $P = \{ f \in L_{\infty [a,b]} : f(x) \geq 0 \mbox{ casi en todas partes } \}   $. 
     \begin{tikzpicture}

        %Función caracteristica 
    % Eje X
    \draw[-] (-1,0) -- (3,0) node[right] {$x$};
    % Eje Y
    \draw[-] (0,-1) -- (0,2) node[above] {$y$};

    % Función característica de [0,1]
    \draw[-, blue] (0,1) -- (1,1);
    \draw[dashed,blue] (1,1) -- (1,0);
    \node at (1, -1) {$\chi_{[0,t_{1}]}(x)$};
    \node at (1, -0.5) {$t_{1}$};

        %Función caracteristica 
    % Eje X
    \draw[-] (5,0) -- (9,0) node[right] {$x$};
    % Eje Y
    \draw[-] (6,-1) -- (6,2) node[above] {$y$};

    % Función característica de [0,2]
    \draw[-, red] (6,1) -- (8,1);
    \draw[dashed,red] (8,1) -- (8,0);
    \node at (8, -1) {$\chi_{[0, t_{2} ]}(x)$};
   \node at (8, -0.5) {$t_{2}$};
\end{tikzpicture}
 \end{block}

\end{frame}
    
%%%%%%%%%%%%%%%%%%%%%%%%%%%%%%%%%%%%%%%%%%%%%%%%%%%%%%%%%%%%%%%%%%%%%%%%%%%%%%%%%%%%%%%%%%%%%%%%%%%%%%%%%%%%%%%%%%%%%%%%%%%%%%%%%%%%%%%%%%%%%%%%%%%%%%%%%%%%%%%%%%%%%%%%%%%%%%%%%%%

\begin{frame}{\textcolor{purple}{No todas las funciones crecientes son de variación acotada o débil de variación acotada}}
    \begin{itemize}
        \item \textbf{Ejemplo:} 
        \begin{block}{}
            Sean $X=C_{\mathbb{R}}^{1}[0,1]$, $||f|| = ||f||_{\infty} + ||f'||_{\infty}$ y $P_0 = \{f \in X: f(0)=0 \  y  \  f(t)\geq 0, \forall t \in[0,1]\}$.\\  $( X, ||\cdot ||, P_0)$ es un ENO.
        \end{block}
        \item Se probará que existe una función creciente que   no es débil de variación acotada y por lo tanto no es de variación de acotada. 
      
    \end{itemize}
\end{frame}




%%%%%%%%%%%%%%%%%%%%%%%%%%%%%%%%%%%%%%%%%%%%%%%%%%%%%%%%%%%%%%%%%%%%%%%%%%%%%%%%%%%%%%%%%%%%%%%%%%%%%%%%%%%%%%%%%%%%%%%%%%%%%%%%%%


\begin{frame}{}
    \begin{itemize}
        \item Función creciente
          \begin{block}{}
            Sea $ f: [0,1] \longrightarrow X  $ dada por 
                 \[ f(t) = \left\{ \begin{array}{lcc}
             \bar{0} , &   si  &  t=1 , \\
             \\ x_{n} ,  &  si  & t \in [1-\frac{1}{n},1-\frac{1}{n+1}) . 
             \end{array}
   \right. \]
   Donde $x_{n}, \bar{0} \in X $ con $x_{n}(s) = - s^{n} $ y $ \bar{0} (s) = 0 $.
        \end{block}
        \item Función lineal y continua
        \begin{block}{}
            $ D: X \longrightarrow \mathbb{R} $ dado por $D(f)= {f}'(1) $.\\
            Es fácil ver que $  |D(f)|= |{f}'(1)| \leq || f' ||_{\infty} \leq || f|| $.
        \end{block}
    \end{itemize}
\end{frame}




\begin{frame}{}
    Por último, 
     \begin{block}{}
                 \[ (D \circ  f)(t) = \left\{ \begin{array}{lcc}
             0 , &   si  &  t=1 , \\
             \\ -n ,  &  si  & t \in [1-\frac{1}{n},1-\frac{1}{n+1}) . 
             \end{array}
   \right. \]
   Ya que  $x_{n}' (t) = - n t^{n-1} $ y $D(x_{n}) = x_{n}' (1) $.
  
        \end{block}

       Sea $m \in \mathbb{N} $ se define $  P_{m} = \{  0,  1 - \frac{1}{m+1}, 1 \}  $ se tiene que 
   \begin{block}{}
      $ V_{a}^{b} \big( (D \circ  f) , P_{m} \big) \geq | (D \circ  f) (1) - (D \circ  f) ( 1 - \frac{1}{m+1}) | = m+1$. 
   \end{block}
   Por lo tanto, $(D \circ  f)$ no es de variación acotada.%, es decir, f no es débilmente de variación acotada. 
  
\end{frame}





%%%%%%%%%%%%%%%%%%%%%%%%%%%%%%%%%%%%%%%%%%%%%%%%%%%%%%%%%%%%%%%%%%%%%%%%%%%%%%%%%%%%%%%%%%%%%%%%%%%%%%%%%%%%%%%%%%%%%%%%%%%%%%%%%%%%%%%%%%%%%%%%%%%%%%%%%%%%%%%%%%%%%%%%%%%%%%%%%%%%%%%%%%%%%%%%%%%%%%%%%%%%%%%%%%%%%%%%%%%%%%%%%%%%%%%%%%%%%%%%%%%%%%%%%%%%%%%%%%%%%%%%%%%%%%%%%%%%%%%%%%%%%%%%%%%%%%%%%%%%%%%%%%%%%%%%%%%%%%%%%%%%%%%%%%%%%%%%%%%%%%%%%%%%%%%%%%%%%%%%%%%%%%%%%%%%%%%%%%%%%%%%

\begin{frame}{}
\textbf{Lema:} Sea $X$ un ENO. Son equivalentes
    \begin{itemize}
         \item $ X_{+} $ es normal.
        \item  Si  toda $ f: [a,b] \longrightarrow X  $ creciente es de variación acotada débil.
        \item Si toda sucesión creciente $ \{ x_{n} \} $ en $X$ es converge, si y sólo si, existe $ \{ x_{n_{k}} \} $ una subsucesión de $ \{ x_{n} \} $ convergente. 
       
    \end{itemize}
\end{frame}

\begin{frame}{}
    \begin{block}{Teorema}
        Si $X$ es un espacio normado ordenado con cono extendible, entonces las funciones crecientes son de variación acotada fuerte. 
    \end{block}
\end{frame}
\section{Descomposición de Jordan}
\begin{frame}{}
    \begin{block}{Definición}
        Un espacio normado ordenado $X$ cumple la descomposición de Jordan si para toda $ f: [a,b] \longrightarrow X $ de variación acotada, existen $  f_{1}, \ f_{2} : [a,b] \longrightarrow X  $ funciones crecientes tales que $ f=f_{1} - f_{2} $.
    \end{block}
    \begin{block}{Teorema}
        Si $X$ cumple la descomposición de Jordan para funciones de variación acotada fuerte, entonces $ X_{+} $ es generador.
    \end{block}
    \begin{block}{Teorema}
        Si  $ X_{+} $ es sólido, entonces $X$ cumple la descomposición de Jordan para funciones de variación acotada fuerte.
    \end{block}
\end{frame}

\begin{frame}{Espacios estudiados}
    \begin{block}{Definición}
        Sea $\left( X, || \cdot  ||, P \right) $ un espacio normado ordenado.\\ Se definen 
		$ BV ( [a,b], X ) = \{   f: [a,b] \longrightarrow X : Var_{a}^{b} (f, X) < \infty \} $\\
		y $ BV([a,b],X)_{0} = \{ f \in BVX : f(a) = 0 \}  $.
    \end{block}
    $  || f ||_{BV} = Var_{a}^{b}  ( f , X ) + ||  f(a) || $ es una norma.
\end{frame}

\begin{frame}{}
    \begin{block}{Lema}
        Sea $X$ un espacio normado ordenado. Entonces $ (  BV([a,b],X) ,  || \cdot ||_{BV} , P_{c} ) $ es un espacio normado ordenado, con el cono $ P_{c} = \{ f \in  BV([a,b],X) : f \mbox{ es creciente y } f \geq 0 \} $.
    \end{block}
   Si $ X_{+} $ es normal, entonces $ P_{c} $ es normal. 
\end{frame}

\begin{frame}{}
    \begin{block}{Teorema}
        Sea $X$ un ENO con cono sólido $P$, $\mu $ una cota superior en orden de la esfera y $ f : [a,b] \longrightarrow X  $ una función vectorial de variación acotada. Entonces las funciones vectoriales $ M_{f}, N_{f} : [a,b] \longrightarrow X $ definidas por 
		\begin{align*}
			M_{f} (x) & = \frac{1}{2} \Bigg[ Var_{a}^{x} (f, X) \mu   + f (x) - f(a) \Bigg] ,\\
			N_{f} (x) & = \frac{1}{2} \Bigg[  Var_{a}^{x} (f, X) \mu   -  f (x) + f(a) \Bigg] 
		\end{align*}
		son funciones crecientes de variación acotada tal que
		\begin{enumerate}
			\item $ f (x) = M_{f} (x) - N_{f} (x) + f(a)  $,
			\item $  Var_{a}^{b} ( M_{f} , X), Var_{a}^{b} ( N_{f} , X) \leq  \frac{1 + || \mu || }{2}   Var_{a}^{b} (f, X) $,
			\item $ Var_{a}^{b} ( M_{f} + N_{f} , X) =  Var_{a}^{b} (f, X) || \mu || $.
			\item $ Var_{a}^{b} ( M_{f} - N_{f} , X) =  Var_{a}^{b} (f, X)  $.
			\item $ || f ||_{BV} = || M_{f} - N_{f} ||_{BV} + || f(a ) || $.
			\item Si además $X$ es un álgebra,entonces  $ BV $ es un álgebra. 
		\end{enumerate}
    \end{block}
\end{frame}



%%%%%%%%%%%%%%%%%%%%%%%%%%%%%%%%%%%%%%%%%%%%%%%%%%%%%%%%%%%%%%%%%%%%%%%%%%%%%%%%%%%%%%%%%%%%%%%%%%%%%%%%%%%%%%%%%%%%%%%%%%%%%%%%%%%%%%%%%%%%%%%%%%%%%%%%%%%%%%%%%%%%%%%%%%%%%%%%%%%%%%%%%%%%%%%%%%%%%%%%%%%%%%%%%%%%%%%%%%%%%%%%%%%%%%%%%%%%%%%%%%%%%%%%%%%%%%%%%%%%%%%%%%%%%%%%%%%%%%%%%%%%%%%%%%%%%%%%%%%%%%%%%%%%%%%%%%%%%%%%

\begin{frame}{Trabajo a futuro}
   


    \begin{block}{Definición:}
        Una función $ f:[a,b] \longrightarrow X $ acotada es \textbf{Riemann-Stieltjes} respecto a $ g:[a,b] \longrightarrow \mathbb{R} $ si existe $ A \in X  $ tal que 
        
        dado $ \epsilon > o  $, existe $ \delta >0  $ de $ [a,b] $ tal que $ \{ (I_{i} , \xi_{i} ) : i = 1,..., n \}   $ es una partición con norma menor $  \delta  $ de $ [a,b]  $, se cumple que
        \begin{equation*}
            || \sum_{i=1}^{n}   f( \xi_{i} ) (g( t_{i} )  - g(t_{i-1}) ) - A  || < \epsilon .
        \end{equation*}
        
    \end{block}
\end{frame}




%%%%%%%%%%%%%%%%%%%%%%%%%%%%%%%%%%%%%%%%%%%%%%%%%%%%%%%%%%%%%%%%%%%%%%%%%%%%%%%%%%%%%%%%%%%%%%%%%%%%%%%%%%%%%%%%%%%%%%%%%%%%%%%%%%%%%%%%%%%%%%%%%%%%%%%%%%%%%%%%%%%%%%%%%%%%%%%%%%%%%%%%%%%%%%%%%%%%%%%%%%%%%%%%%%%%%%%%%%%%%%%%%%%%%%%%%%%%%%%%%%%%%%%%%%%%%%%%%%%%%%%%%%%%%%%%%%%%%%%%%%%%%%%%%%%%%%%%%%%%









	


%%%%%%%%%%%%%%%%%%%%%%%%%%%%%%%%%%%%%%%%%%%%%%%%%%%%%%%%%%%%%%%%%%%%%%%%%%%%%%%%%%%%%%%%%%%%%%%%%%%%%%%%%%%%%%%%%%%%%%%%%%%%%%%%%%%%%%%%%%%%%%%%%%%%%%%%%%%%%%%%%%%%%%%%%%%%%%%%%%%%%%%%%%%%%%%%%%%%%%%%%%%%%%%%%%%%%%%%%%%%%%%%%%%%%%%%%%%%%%%%%%%%%%%%%%%%%%%%%%%%%%%%%%%%%%%%%%%%%%%%%%%%%%%%%%%%%%%%%%%%



%%%%%%%%%%%%%%%%%%%%%%%%%%%%%%%%%%%%%%%%%%%%%%%%%%%%%%%%%%%%%%%%%%%%%%%%%%%%%%%%%%%%%%%%%%%%%%%%%%%%%%%%%%%%%%%%%%%%%%%%%%%%%%%%%%%%%%%%%%%%%%%%%%%%%%%%%%%%%%%%%%%%%%%%%%%%%%%%%%%%%%%%%%%%%%%%%%%%%%%%%%%%%%%%%%%%%%%%%%%%%%%%%%%%%%%%%%%%%%%%%%%%%%%%%%%%%%%%%%%%%%%%%%%%%%%%%%%%%%%%%%%%%%%%%%%%%%%%%%%%%%%%%%%%%%%%%%%%%%%%%%%%%%%%%%%%%




%\section{Descomposición de Jordan}

%   %%%%%%%%%%%%%%%%%%%%%%%%%%%%%%%%%%%%%%%%%%%%%%%%%%%%%%%%%%%%%%%%%%%%%%%%%%%%%%%%%%%%%%%%%%%%%%%%%%%%%%%%%%%%%%%%%%%%%%%%%%%%%%%%%%%%%%%%%%%%%%%%%%%%%%%%%%%%%%%%%%%%%%%%%%%%%%%%%%%%%%%%%%%%%%%%%%%%%%%%%%%%%%%%%%%%%%%%%%%%%%%%%%%%%%%%%%%%%%%%%%%%%%%%%%%%%%%%%%%%%%%%%%%%%%%%%%%%%%%%%%%%%%%%%%%%%%%%%%%%%%%%%%%%%%%%%%%%%%%%%%%%%%%%%%%%%%

\begin{frame}{}
    \begin{center}
        \Huge  Algunos resultados sobre la ``Descomposición de Jordan"  
    \end{center}
\end{frame}


%%%%%%%%%%%%%%%%%%%%%%%%%%%%%%%%%%%%%%%%%%%%%%%%%%%%%%%%%%%%%%%%%%%%%%%%%%%%%%%%%%%%%%%%%%%%%%%%%%%%%%%%%%%%%%%%%%%%%%%%%%%%%%%%%%%%%%%%%%%%%%%%%%%%%%%%%%%%%%%%%%%%%%%%%%%%%%%%%%%%%%%%%%%%%%%%%%%%%%%%%%%%%%%%%%%%%%%%%%%%%%%%%%%%%%%%%%%%%%%%%%%%%%%%%%%%%%%%%%%%%%%%%%%%%%%%%%%%%%%%%%%%%%%%%%%%%%%%%%%%

\begin{frame}{Condiciones necesarias para  la descomposición de Jordan}
\begin{itemize}
    \item \textbf{Definición:}
      \begin{block}{}
          Se dice que un espacio normado ordenado $X$ satisface la descomposición de Jordan si toda función de variación acotada $ g,h:[a,b]\to X$, existen $f:[a,b]\to X$ crecientes tales que $ f=g-h $.
      \end{block}
    \item \textbf{Lema}
\end{itemize}
    \begin{block}{}
        Si $f:[a,b]\to X$ cumple la descomposición de Jordan, entonces $f$ es orden acotado. 
    \end{block}
\end{frame}


%%%%%%%%%%%%%%%%%%%%%%%%%%%%%%%%%%%%%%%%%%%%%%%%%%%%%%%%%%%%%%%%%%%%%%%%%%%%%%%%%%%%%%%%%%%%%%%%%%%%%%%%%%%%%%%%%%%%%%%%%%%%%%%%%%%%%%%%%%%%%%%%%%%%%%%%%%%%%%%%%%%%%%%%%%%%%%%%%%%%%%%%%%%%%%%%%%%%%%%%%%%%%%%%%%%%%%%%%%%%%%%%%%%%%%%%%%%%%%%%%%%%%%%%%%%%%%%%%%%%%%%%%%%%%%%%%%%%%%%%%%%%%%%%%%%%%%%%%%%%

    \begin{frame}{}
        \begin{block}{Teorema}
            Sea $X$ un ENO. Si $P$ no es generador, entonces existe  $f : [a,b] \longrightarrow X $  de variación acotada y que no cumple la descomposición de Jordan.
        \end{block}
        Demostración: Sea $ x_{0} \in X  $ tal que $ x_{0} \notin P - P $. Se define $ f: [a,b] \longrightarrow X $ 
        
        \[ f(t) = \left\{ \begin{array}{lcc}
             2 x_{0} , &   si  &  t=b , \\
             \\  x_{0} , &  si & t=a , \\
             \\ 0 ,  &  si  & t \in (a,b) . 
             \end{array}
   \right. \]
        Si existieran $g,h: [a,b] \longrightarrow X$ crecientes tal que $ f = g - h $ entonces 
        $ x_{0} = f(b) - f(a) = (g(b) - g(a) ) - ( h(b) - h(a)) . $
    \end{frame}

%%%%%%%%%%%%%%%%%%%%%%%%%%%%%%%%%%%%%%%%%%%%%%%%%%%%%%%%%%%%%%%%%%%%%%%%%%%%%%%%%%%%%%%%%%%%%%%%%%%%%%%%%%%%%%%%%%%%%%%%%%%%%%%%%%%%%%%%%%%%%%%%%%%%%%%%%%%%%%%%%%%%%%%%%%%%%%%%%%%%%%%%%%%%%%%%%%%%%%%%%%%%%%%%%%%%%%%%%%%%%%%%%%%%%%%%%%%%%%%%%%%%%%%%%%%%%%%%%%%%%%%%%%%%%%%%%%%%%%%%%%%%%%%%%%%%%%%%%%%%%%%%%%%%%

\begin{frame}{}
    \begin{itemize}
        \item \textbf{Ejemplo}
        \begin{block}{}
            	El espacio $(L_1[0,1], \|\cdot\|_1)$ con el orden usual de funciones. El cono $L_1[0,1]_+$ es generador ya que toda función en $L_1[0,1]$ es la diferencia de su parte positiva  parte negativa. Ver el ejemplo 2.5.1 en \cite{Gou2004}. %  shows $L_1[0,1]_+$ is an extensible cone.
        \end{block}
        \item \textbf{Función}
        \begin{block}{}
            Sea $f:[0,1] \longrightarrow L_1 [0,1]$  definida por  $f(0)=0$ (función nula) y para todo  $t \in (0,1]$, $f(t)=T_t$, donde $T_t:[0,1] \longrightarrow \mathbb{R}$ está dada por

	$$ T _t(s)=\left\{
	\begin{array}{ll}
		n_0 , & \hbox{$s\in[0,\frac{1}{(n_0 )^3}]$;} \\
		0, & \hbox{otherwise,}
	\end{array}
	\right.
	$$
donde $n_0 = [\frac{1}{t}]$.
        \end{block}
    \end{itemize}
\end{frame}


%%%%%%%%%%%%%%%%%%%%%%%%%%%%%%%%%%%%%%%%%%%%%%%%%%%%%%%%%%%%%%%%%%%%%%%%%%%%%%%%%%%%%%%%%%%%%%%%%%%%%%%%%%%%%%%%%%%%%%%%%%%%%%%%%%%%%%%%%%%%%%%%%%%%%%%%%%%%%%%%%%%%%%%%%%%%%%%%%%%%%%%%%%%%%%%%%%%%%%%%%%%%%%%%%%%%%%%%%%%%%%%%%%%%%%%%%%%%%%%%%%%%%%%%%%%%%%%%%%%%%%%%%%%%%%%%%%%%%%%%%%%%%%%%%%%%%%%%%%%%%%%%%%%%%

\begin{frame}{}
    	Let $P= \{t_0 =0, t_1, ... , t_{k}=1 \} \in \mathcal{P}[a,b] $ and $m_{i}= [\frac{1}{t_{i}}] \in \mathbb{N}  $ for all $i=1,...k.$ 
	Since $1= m_{n} \leq m_{n-1} \leq \cdots \leq m_{1}$, then 
	$ \frac{1}{(m_1)^3} \leq \frac{1}{(m_2)^3} \leq  \cdots \leq \frac{1}{(m_n)^3}=1  $.\\
	Define a set  $ \{ t'_{0}, t'_{1}, ... , t'_{n}  \}$ by $ t'_{0}=0, t'_{1}= t_{1} $ and $ t'_{j}= \min \{ t_{i} \in P- \{ 0 \}  :  \big[ \frac{1}{t_{i}}\big]  <  \big[ \frac{1}{t'_{j-1}}\big]  \} $, then  $ \{ t'_{0}, t'_{1}, ... , t'_{n}  \} \subset P $.\\ 
Note that if $ m_{i} =\big[ \frac{1}{(m_{i-1})3} \big] =\big[ \frac{1}{(m_{i})3} \big] = m_{i+1} $, then $ \int_{0}^{1} |T_{t_{i}}(s) - T_{t_{i-1}}(s) | ds= 0 $.
\end{frame}

%%%%%%%%%%%%%%%%%%%%%%%%%%%%%%%%%%%%%%%%%%%%%%%%%%%%%%%%%%%%%%%%%%%%%%%%%%%%%%%%%%%%%%%%%%%%%%%%%%%%%%%%%%%%%%%%%%%%%%%%%%%%%%%%%%%%%%%%%%%%%%%%%%%%%%%%%%%%%%%%%%%%%%%%%%%%%%%%%%%%%%%%%%%%%%%%%%%%%%%%%%%%%%%%%%%%%%%%%%%%%%%%%%%%%%%%%%%%%%%%%%%%%%%%%%%%%%%%%%%%%%%%%%%%%%%%%%%%%%%%%%%%%%%%%%%%%%%%%%%%%%%%%%%%%

\begin{frame}{}
    	It follows that 
	\begin{align*}
		\sum_{i=1}^{k} || f(t_{i-1}) - f(t_{i})||_{1} &= \sum_{i=1}^{k} \int_{0}^{1} |T_{t_{i-1}}(s) - T_{t_{i}}(s) |ds \\
		&= \sum_{i=1}^{n} \int_{0}^{1} |T_{t'_{i}}(s) - T_{t'_{i-1}}(s) |ds \\
		&= \int_{0}^{\frac{1}{(m'_1)3}} m'_1 ds \\
  & + \sum_{i=2}^{n} \left[ \int_{0}^{\frac{1}{(m'_{i-1})3}} ( m'_{i-1} - m'_{i} ) ds+ \int_{\frac{1}{(m'_{i-1})3}}^{\frac{1}{(m'_i)3}} m'_ids  \right] 
	\end{align*} 

\end{frame}


%%%%%%%%%%%%%%%%%%%%%%%%%%%%%%%%%%%%%%%%%%%%%%%%%%%%%%%%%%%%%%%%%%%%%%%%%%%%%%%%%%%%%%%%%%%%%%%%%%%%%%%%%%%%%%%%%%%%%%%%%%%%%%%%%%%%%%%%%%%%%%%%%%%%%%%%%%%%%%%%%%%%%%%%%%%%%%%%%%%%%%%%%%%%%%%%%%%%%%%%%%%%%%%%%%%%%%%%%%%%%%%%%%%%%%%%%%%%%%%%%%%%%%%%%%%%%%%%%%%%%%%%%%%%%%%%%%%%%%%%%%%%%%%%%%%%%%%%%%%%%%%%%%%%%

\begin{frame}{}
    	
	\begin{align*}
		&= \sum_{i=1}^{n} \int_{0}^{1} |T_{t'_{i}}(s) - T_{t'_{i-1}}(s) |ds \\
		&= \int_{0}^{\frac{1}{(m'_1)3}} m'_1 ds+ \sum_{i=2}^{n} \left[ \int_{0}^{\frac{1}{(m'_{i-1})3}} ( m'_{i-1} - m'_{i} ) ds+ \int_{\frac{1}{(m'_{i-1})3}}^{\frac{1}{(m'_i)3}} m'_ids  \right] \\
		&\leq  \int_{0}^{\frac{1}{(m'_1)3}} m'_1ds + \sum_{i=2}^{n} \left[ \int_{0}^{\frac{1}{(m'_{i-1})3}}  m'_{i-1}ds  + \int_{0}^{\frac{1}{m'_{i}3}} m'_ids  \right]  \\
		& = \frac{1}{(m'_1)^2} + \sum_{i=2}^{n} \left[ \frac{1}{(m'_{i-1})^2} +  \frac{1}{(m'_{i})^2} \right] \\
		& \leq  \sum_{i=1}^{n-1} \frac{1}{(m'_{i})^2} + \sum_{i=1}^{n} \frac{1}{(m'_{i})^2} \leq \frac{\pi^2}{6} + \frac{\pi^2}{6} = \frac{\pi^2}{3} .
	\end{align*} 
Por lo tanto, $f$ es de variación acotada, pero no acotada en orden. Así, $f$  no cumple la descomposición de Jordan.
\end{frame}

%%%%%%%%%%%%%%%%%%%%%%%%%%%%%%%%%%%%%%%%%%%%%%%%%%%%%%%%%%%%%%%%%%%%%%%%%%%%%%%%%%%%%%%%%%%%%%%%%%%%%%%%%%%%%%%%%%%%%%%%%%%%%%%%%%%%%%%%%%%%%%%%%%%%%%%%%%%%%%%%%%%%%%%%%%%%%%%%%%%%%%%%%%%%%%%%%%%%%%%%%%%%%%%%%%%%%%%%%%%%%%%%%%%%%%%%%%%%%%%%%%%%%%%%%%%%%%%%%%%%%%%%%%%%%%%%%%%%%%%%%%%%%%%%%%%%%%%%%%%%%%%%%%%%%
\begin{frame}{Condiciones suficientes para la descomposición de Jordan }
\begin{block}{Definición:}
    Sean $(X, || \cdot || , P)  $ un ENO, $ \mu  \in P $ y la esfera $ S(X) = \{ x \in X : || x || =1 \}   $.  $S(X)$  es acotada superiormente  en orden por $ \mu  $ si y sólo si\\  $ x \leq || x || \mu  $  para cada $ x \in X $.
\end{block}
\begin{block}{Lema}
    Sean $X  $ un ENO y  $ \mu   $ una cota superior  de $S(X)$. Entonces toda  $ f :[a,b] \longrightarrow X $ de variación acotada cumple la descomposición de Jordan.
\end{block}

\end{frame}


%%%%%%%%%%%%%%%%%%%%%%%%%%%%%%%%%%%%%%%%%%%%%%%%%%%%%%%%%%%%%%%%%%%%%%%%%%%%%%%%%%%%%%%%%%%%%%%%%%%%%%%%%%%%%%%%%%%%%%%%%%%%%%%%%%%%%%%%%%%%%%%%%%%%%%%%%%%%%%%%%%%%%%%%%%%%%%%%%%%%%%%%%%%%%%%%%%%%%%%%%%%%%%%%%%%%%%%%%%%%%%%%%%%%%%%%%%%%%%%%%%%%%%%%%%%%%%%%%%%%%%%%%%%%%%%%%%%%%%%%%%%%%%%%%%%%%%%%%%%%%%%%%%%%%
\begin{frame}{ }

\begin{block}{Lema}
    Sean $X  $ un ENO y  $ \mu   $ una cota superior  de $S(X)$. Entonces toda  $ f :[a,b] \longrightarrow X $ de variación acotada cumple la descomposición de Jordan.

\end{block}
   Demostración:\\
   Sean  $ g,h :[a,b] \longrightarrow X $ dadas por $ g(t) = Var_{a}^{t} (f, X) \mu - f(t) $ y $ h(t) = Var_{a}^{t} (f, X) \mu $.\\
 Como  $ \mu   $ una cota superior  de $S(X)$, entonces $f$, $g$ son crecientes y \\
 $  f= h - g $.
\end{frame}



%%%%%%%%%%%%%%%%%%%%%%%%%%%%%%%%%%%%%%%%%%%%%%%%%%%%%%%%%%%%%%%%%%%%%%%%%%%%%%%%%%%%%%%%%%%%%%%%%%%%%%%%%%%%%%%%%%%%%%%%%%%%%%%%%%%%%%%%%%%%%%%%%%%%%%%%%%%%%%%%%%%%%%%%%%%%%%%%%%%%%%%%%%%%%%%%%%%%%%%%%%%%%%%%%%%%%%%%%%%%%%%%%%%%%%%%%%%%%%%%%%%%%%%%%%%%%%%%%%%%%%%%%%%%%%%%%%%%%%%%%%%%%%%%%%%%%%%%%%%%%%%%%%%%%


\begin{frame}{ }
\textbf{Lema:}% Sea $(X, || \cdot || , P)  $ un ENO y la esfera $ S(X) = \{ x \in X : || x || =1 \}   $.
	\begin{itemize}

	\item   La esfera  es acotada superiormente  en orden si y sólo si  $  P $ es sólido. 
		   
 
	\begin{tikzpicture}
	
          

		%Dibuja los ejes figura 1
		\draw[->] (-1.5,0) -- (2.5,0);
		\draw [->] (0,-1.5) -- (0,2.2);
		%Circulo de figura 1 
		 \draw[color=blue] (0,0) circle (0.5cm);
		
			%Dibuja los ejes figura 2
		\draw[->] (3.5,0) -- (7.5,0);
		\draw [->] (5,-1.5) -- (5,2.2);
		
			%%%Crea el cono P en la figura 2 %%%
		\fill[fill=yellow]
		(5,0)
		-- (8.5,2.7) 
		-- (5.3,3.3) node[fill=red!20,draw,double,rounded corners] {$ P $};
		%Circulo en la figura 2
		\draw[color=blue] (6,2) circle (0.5cm);
		 \node[ color=blue] (2) at (6,2) { $ \cdot  \mu  $};
		%Linea verde que une figurta 1 y figura 2 
      \draw[<->,color=green] (0.2,0.2) to[bend left] (6,1.8);
      
      
      
      \node[draw, color=blue] (2) at (-1,2.6) { $|| x || < r $};
      \node[draw, color=green] (3) at (0.5,2.6) { $ \iff  $};
      \node[draw, color=blue] (4) at (3,2.6) { $ || \mu - (\mu - x) || < r $};
   
\end{tikzpicture}
	
	
	
	\end{itemize}
\end{frame}






%%%%%%%%%%%%%%%%%%%%%%%%%%%%%%%%%%%%%%%%%%%%%%%%%%%%%%%%%%%%%%%%%%%%%%%%%%%%%%%%%%%%%%%%%%%%%%%%%%%%%%%%%%%%%%%%%%%%%%%%%%%%%%%%%%%%%%%%%%%%%%%%%%%%%%%%%%%%%%%%%%%%%%%%%%%%%%%%%%%%%%%%%%%%%%%%%%%%%%%%%%%%%%%%%%%%%%%%%%%%%%%%%%%%%%%%%%%%%%%%%%%%%%%%%%%%%%%%%%%%%%%%%%%%%%%%%%%%%%%%%%%%%%%%%%%%%%%%%%%%%%%%%%%%%

\begin{frame}{}
    \begin{itemize}
        \item \textbf{Ejemplo:} Un espacio sin esfera acotada que cumple la descomposición de Jordan. 
        \begin{block}{}
            	$ (l^{1}, || \cdot ||_{1} ) $ con $l^{1}_{+} = \{ \{a_{n} \} \in l^{1} : a_{n} \geq 0 \mbox{ para todo } n \in \mathbb{N} \} $ cumple la descomposición de Jordan.
        \end{block}
        	Sea $f:[a,b]\rightarrow l^{1}$ uy $f_{j} :[a,b]\rightarrow \mathbb{R} $ con $ j \in \mathbb{N} $ tal que para  $t\in [a,b]$
		\[f(t)=\sum_{j=1}^{\infty} f_{j}(t)e_{j} , \]
		donde $ e_{j} $ es un vector unitario sobre
	 $l^{1}$.\\	
	 Si  $f$ es v. a., entonces $ f_{j} $ es v. a. y $ f_{j} (x) = V (f_{j}, [a,x] ) - [ V (f_{j}, [a,x] ) - f_{j} (x) ] $. 
	
	
    \end{itemize}
\end{frame}

\begin{frame}{}
    $ \{ V (f_{j}, [a,x] ) \} \in l^{1} $ ya que  $ 0 \leq  V (f_{j}, [a,x] ) \leq  V (f_{j}, [a,b] ) $ y así $ \sum_{j=1}^{\infty} V (f_{j}, [a,x] ) \leq \sum_{j=1}^{\infty}  V (f_{j}, [a,b] ) \leq V (f, [a,b] ) $.
	
	Por tanto,  $$ f_{2}(x) = \{ V (f_{j}, [a,x] ) \} , \  f_{1} (x) = \{ V (f_{j}, [a,x] ) - f_{j} (x) \} \in l^{1} $$
 y así
 $f= f_{2} - f_{1}$.
\end{frame}



%%% ─────────────────────────────────────────────────────────────────────────────
%%  El espacio BV([a,b],X) como ENO
%% ─────────────────────────────────────────────────────────────────────────────
\section{El espacio $\BVX$}

\begin{frame}{}
\begin{block}{}
\centering\large
El espacio $BV([a,b],X)$
\end{block}
\end{frame}

%──────────────────────────────────────────────────────────────────────────────

\begin{frame}{Definición y estructura de Banach}
\framesubtitle{$BV([a,b],X)$ hereda la completitud de $X$}

\begin{definition}
Sea $(X,\|\cdot\|,\Xp)$ un ENO. Se define
\[
\BVX = \bigl\{f:[a,b]\to X : \Var_a^b(f,X)<\infty\bigr\}
\]
con norma $\|f\|_{\BVX} = \|f(a)\| + \Var_a^b(f,X)$.
\end{definition}

\medskip

\begin{theorem}
\begin{enumerate}
  \item Para toda $f\in\BVX$,
        $\|f\|_\infty \le \|f\|_{\BVX}$.
  \item $(\BVX,\|\cdot\|_{\BVX})$ es un \textbf{espacio de Banach}.
\end{enumerate}
\end{theorem}

\end{frame}

%──────────────────────────────────────────────────────────────────────────────

\begin{frame}{$\BVX$ como espacio normado ordenado}
\framesubtitle{El orden se hereda del cono de funciones crecientes}

Sea $P_c = \{f\in\BVX : f \text{ creciente},\ f\ge 0\}$.

\medskip

\begin{theorem}[\cite{perez2022jordan}]
\begin{itemize}
  \item[--] $P_c$ es un cono para $(\BVX,\|\cdot\|_{\BVX})$.
  \item[--] Si $\Xp$ es \textbf{sólido}, entonces $P_c$ es cono
            \textbf{generador} de $\BVX$.
  \item[--] Si $P_c$ es cono generador de $\BVX$, entonces $\Xp$
            es generador.
\end{itemize}
\end{theorem}

\end{frame}

%──────────────────────────────────────────────────────────────────────────────

\begin{frame}{Propiedades del cono de $\BVX$}
\framesubtitle{Normalidad y transferencia de estructura}

\begin{block}{Normalidad}
Si $\Xp$ es \textbf{normal}, entonces $P_c$ es cono \textbf{normal}
en $\BVX$.
\end{block}

\medskip

\begin{alertblock}{Lectura}
Las propiedades del cono se transfieren: $\BVX$ es un ENO y
sus propiedades reflejan las del cono de $X$.
\end{alertblock}

\end{frame}

%──────────────────────────────────────────────────────────────────────────────

\begin{frame}{Estructura algebraica de $\BVX$}
\framesubtitle{Cuando $X$ es álgebra}

\begin{theorem}[\cite{perez2022jordan}]
Si $X$ es un álgebra con $\|fg\|\le\|f\|\,\|g\|$,
entonces $(\BVX,3\|\cdot\|_{\BVX})$ también es un \textbf{álgebra}.
\end{theorem}

\medskip

\begin{alertblock}{Importancia}
Esta estructura permitirá aplicar teoremas de punto fijo en $\BVX$
como espacio normado ordenado con cono y álgebra.
\end{alertblock}

% TODO: agregar referencia exacta para el álgebra BV.

\end{frame}


%%%%%%%%%%%%%%%%%%%%%%%%%%%%%%%%%%%%%%%%%%%%%%%%%%%%%%%%%%%%%%%%%%%%%%%%%%%%%%%%%%%%%%%%%%%%%%%%%%%%%%%%%%%%%%%%%%%%%%%%%%%%%%%%%%%%%%%%%%%%%%%%%%%%%%%%%%%%%%%%%%%%%%%%%%%%%%%%%%%%%%%%%%%%%%%%%%%%%%%%%%%%%%%%%%%%%%%%%%%%%%%%%%%%%%%%%%%%%%%%%%%%%%%%%%%%%%%%%%%%%%%%%%%%%%%%%%%%%%%%%%%%%%%%%%%%%%%%%%%%%%%%%%%%%







%%%%%%%%%%%%%%%%%%%%%%%%%%%%%%%%%%%%%%%%%%%%%%%%%%%%%%%%%%%%%%%%%%%%%%%%%%%%%%%%%%%%%%%%%%%%%%%%%%%%%%%%%%%%%%%%%%%%%%%%%%%%%%%%%%%%%%%%%%%%%%%%%%%%%%%%%%%%%%%%%%%%%%%%%%%%%%%%%%%%%%%%%%%%%%%%%%%%%%%%%%%%%%%%%%%%%%%%%%%%%%%%%%%%%%%%%%%%%%%%%%%%%%%%%%%%%%%%%%%%%%%%%%%%%%%%%%%%%%%%%%%%%%%%%
%\section{\textcolor{blue}{Conclusiones}}


\begin{frame}{\textcolor{blue}{Conclusiones}}
  \begin{itemize}
      \item \textcolor{green}{Logros}
      \begin{enumerate}[(1)]
    \item  Se dio un ejemplo de que no toda función vectorial $ f: [a,b] \longrightarrow X  $ creciente es débil de variación acotada.
    \item  Se probo que para toda función vectorial $ f: [a,b] \longrightarrow X  $ creciente sea débil de variación acotada, es equivalente a que el cono sea de normal.
   % \item  Se probo que un cono sea sólido es condición suficiente, pero no necesaria para que $X$ cumpla la descomposición de Jordan dando un ejemplo.
    \end{enumerate}
    \item \textcolor{purple}{Perspectivas}
    \begin{itemize}
        \item[(1)] Ver la relación que el cono sea normal a que toda función vectorial sea integrable. 
        \item[(2)]  Ver si es equivalente que el cono extendible a que toda función vectorial sea  de variación acotada. 
        %\item[(3)] Analizar más a profundidad el espacio $ BV([a,b],X) $ como ENO  y aplicar los teoremas de punto fijo que tenemos.       
    \end{itemize}
  \end{itemize}
    
     \end{frame}



%%%%%%%%%%%%%%%%%%%%%%%%%%%%%%%%%%%%%%%%%%%%%%%%%%%%%%%%%%%%%%%%%%%%%%%%%%%%%%%%%%%%%%%%%%%%%%%%%%%%%%%%%%%%%%%%%%%%%%%%%%%%%%%%%%%%%%%%%%%%%%%%%%%%%%%%%%%%%%%%%%%  BIBLIOGRAFÍA  %%%%%%%%%%%%%%%%%%%%%%%%%% %%%%%%%%%%%%%%%%%%%%%%%%%%%%%%%%%%%%%%%%%%%%%%%%%%%%%%%%%%%%%%%%%%%%%%%%%%%%%%%%%%%%%%%%%%%%%%%%%%%%%%%%%%%%%%%%%%%%%%%%%%%%%%%%

\section{\textcolor{blue}{Bibliografía}}
\thispagestyle{empty}
\newpage

\begin{frame}[t,fragile]{Bibliografía}
\begin{thebibliography}{99}
    \bibitem{Baz} \textsc{Agarwal Ravi P., Meehan Maria } y \textsc{O'Regan Donal},	\textit{Fixed point theory and applications}, volumen 141.
	Cambridge University Press. First Edition. 2001.

	\bibitem{cone2007}    \textsc{Aliprantis, Charalambos D} y \textsc{Tourky, Rabee}.  \textit{Cones and duality}. American Mathematical Soc., volumen 84.  \hypertarget{55}{(2007)}.\hyperlink{01}{\beamergotobutton{Regreso}}
	
	 \bibitem{Gou2004} \textsc{Gou D., Yeol J. C. } y \textsc{Jiang Z.}	\textit{Partial Ordering Methods in Nonlinear Problems}, 
	Nova Science Plublishers Inc. Segunda edición.\hypertarget{05}{ (2004)}.
	\hyperlink{22}{\beamergotobutton{Regreso}}

	
 \bibitem{Granas1994} \textsc{Granas, Andrzej} 	\textit{Continuation Method for Contractive Maps}, 
	Journal of the Juliusz Schauder Center. Volume 3. (1994)
	
	
	





\end{thebibliography}
\end{frame}



\thispagestyle{empty}
\newpage

\begin{thebibliography}{99}    
    \bibitem{G.E} \textsc{ HARDY, G.E.}, \textsc{ ROGERS, T.D.}
	 \textit{A generalization of a fixed point theorem of reich}. Canad. Math. Bull, volume 16. (1973).

	\bibitem{OL2016}   \textsc{Jiang, Shujun} y \textsc{Li, Zhilong}.  Article. \textit{Fixed point theorems of order-Lipschitz mappings in Banach algebras}, Fixed Point Theory and Applications.  \hypertarget{03}{(2016)}.
		\hyperlink{33}{\beamergotobutton{Regreso}}
	
	\bibitem{OL2018}    \textsc{Jiang, Shujun} y \textsc{Li, Zhilong}.  Article. \textit{Order-Lipschitz Mappings Restricted with Linear Bounded Mappings in Normed Vector Spaces without Normalities of Involving Cones via Methods of Upper and Lower Solutions}, Faculty of Sciences and Mathematics University of Ni\v{s}. \hypertarget{02}{(2018)}. \hyperlink{11}{\beamergotobutton{Regreso}}



\bibitem{Roque}   \textsc{Luciano Gerardo Roque Vidal,} \textit{TEOREMAS DEL PUNTO FIJO PARA FUNCIONES MONÓTONAS Y SUS APLICACIONES} (Tesis de Licenciatura). FCFM-BUAP. 2018.




\bibitem{Rhoad} \textsc{Rhoades Billy E.} 
	\textit{A comparison of various definitions of contractive mappings}. Transactions of the American Mathematical Society,  volumen 226. (1977).

 \bibitem{ReichNew} \textsc{Shaoyuan Xu}, \textsc{Suyu Cheng} y \textsc{Zuoling Zhou} 
	\textit{Reich´s iterated function systems and well-posedness via fixed point theory}. Fixed Point Theory and Applications, volumen 2015. (2015).

\bibitem{anel}   \textsc{Vázquez Martínez Anel,} \textit{Teorema del Punto Fijo en Espacios de Bananch Ordenados y
		Aplicaciones} (Tesis de Maestría). FCFM-BUAP. 2018.

	\bibitem{perez2022jordan} \textsc{P{\'e}rez-Becerra, Tom{\'a}s and Escamilla-Reyna, Juan Alberto and Luciano-Gerardo, Roque Vidal and Mendoza-Torres, Francisco J and Rodr{\'\i}guez-Tzompantzi, Daniela }  \textit{The Jordan Decomposition in Ordered Normed Spaces}.
   Mediterranean Journal of Mathematics,
  volume 19,
  number 5, Springer,
  2022.
  

	
\end{thebibliography}

\thispagestyle{empty}

%%%%%%%%%%%%%%%%%%%%%%%%%%%%%%%%%%%%%%%%%%%%%%%%%%%%%%%%%%%%%%%%%%%%%%%%%%%%%%%%%%%%%%%%%%%%%%%%%%%%%%%%%%%%%%%%%%%%%%%%%%%%%%%%%%%%%%%%%%%%%%%%%%%%%%%%%%%%%%%%%%%  PARTE DESPUÉS DE LA BIBLIOGRAFÍA %%%%%%%%%%%%%%  %%%%%%%%%%%%%%%%%%%%%%%%%%%%%%%%%%%%%%%%%%%%%%%%%%%%%%%%%%%%%%%%%%%%%%%%%%%%%%%%%%%%%%%%%%%%%%%%%%%%%%%%%%%%%%%%%%%%%%%%%%%%%%%%
\nonumber
\newpage


	\begin{center}
		\Huge ¡Gracias por su atención!	
		
	\end{center}
   %  \begin{center}
    % \Large ¿Preguntas?	
    % \end{center}
	\thispagestyle{empty}
\newpage%
	
%%%%%%%%%%%%%%%%%%%%%%%%%%%%%%%%%%%%%%%%%%%%%%%%%%%%%%%%%%%%%%%%%%%%%%%%%%%%%%%%%%%%%%%%%%%%%%%%%%%%%%%%%%%%%%%%%%%%%%%%%%%%%%%%%%%%%%%%%%%%%%%%%%%%%%%%%%%%%%%%%%%%%%%%%%%%%%%%%%%%%%%%%%%%%%%%%%%%%%%%%%%%%%%%%%%%%%%%%%%%%%%%%%%%%%%%%%%%%%%%%%%%%%%%%%%%%%%%%%%%%%%%%%%%%%
\textbf{Theorem} 		Let $f:[a,b]\rightarrow l^{1}$ be a vector function and $f_{j} :[a,b]\rightarrow \mathbb{R} $ with $ j \in \mathbb{N} $ such that for all  $t\in [a,b]$
		\[f(t)=\sum_{j=1}^{\infty} f_{j}(t)e_{j} , \]
		where $ e_{j} $ is a unit vector on
	 $l^{1}$.\\	
	 If $f$ is bounded variation, then $  f_{j}$ is bounded variation for all  $ j \in \mathbb{N} $  and 
		
		$  V(f,[a,b]) = \sum_{j=1}^{\infty} V(f_{j},[a,b])\ $.\\ 
	Proof: \\ 
	 Let  $P=\{a=t_{0}<\cdots <t_{n}=b\}$ a partition of  $[a,b]$.
		Then
		\[\sum_{i=1}^{n}|f_{j}(t_{i})-f_{j}(t_{i-1}))|\leq \sum_{i=1}^{n}\|f(t_{i})-f(t_{i-1})\|_{1}\leq V(f,[a,b]).\]

\thispagestyle{empty}
\newpage%
	
%%%%%%%%%%%%%%%%%%%%%%%%%%%%%%%%%%%%%%%%%%%%%%%%%%%%%%%%%%%%%%%%%%%%%%%%%%%%%%%%%%%%%%%%%%%%%%%%%%%%%%%%%%%%%%%%%%%%%%%%%%%%%%%%%%%%%%%%%%%%%%%%%%%%%%%%%%%%%%%%%%%%%%%%%%%%%%%%%%%%%%%%%%%%%%%%%%%%%%%%%%%%%%%%%%%%%%%%%%%%%%%%%%%%%%%%%%%%%%%%%%%%%%%%%%%%%%%%%%%%%%%%%%%%%%

  
		Thus, $f_{j}$is bounded variation on $[a,b]$ and
		\[V(f_{j},[a,b])\leq V(f,[a,b]).\]
		
		First, we will prove 	$ \sum_{j=1}^{\infty} V(f_{j},[a,b])\leq V(f,[a,b])$.
		
		
		Next,  let's suppose $ V(f,[a,b]) <  \sum_{j=1}^{\infty} V(f_{j},[a,b])   $, then exists $n_{0} \in \mathbb{N}  $ such that $ V(f,[a,b]) <  \sum_{j=1}^{n_{0}} V(f_{j},[a,b])   $.


Let $ \epsilon = \dfrac{\sum_{j=1}^{n_{0}} V(f_{j},[a,b]) - V(f,[a,b])}{n_{0}} > 0$, then for all $ j \in \{ 1,  \cdots , n_{0}  \} $	 exists $ P_{j} = \{a=x_{0}<\cdots < x_{n_{j}}= b \}  $	a partition of  $[a,b]$ such that 
\begin{equation*}
    \sum_{k=1}^{n_{j}} |f_{j}(x_{k})-f_{j}(x_{k-1}) | > V(f_{j},[a,b]) - \epsilon .
\end{equation*}

\thispagestyle{empty}
\newpage%
	
%%%%%%%%%%%%%%%%%%%%%%%%%%%%%%%%%%%%%%%%%%%%%%%%%%%%%%%%%%%%%%%%%%%%%%%%%%%%%%%%%%%%%%%%%%%%%%%%%%%%%%%%%%%%%%%%%%%%%%%%%%%%%%%%%%%%%%%%%%%%%%%%%%%%%%%%%%%%%%%%%%%%%%%%%%%%%%%%%%%%%%%%%%%%%%%%%%%%%%%%%%%%%%%%%%%%%%%%%%%%%%%%%%%%%%%%%%%%%%%%%%%%%%%%%%%%%%%%%%%%%%%%%%%%%%


	Let $ P = \bigcup_{j=1}^{n_{0}} P_{j} $ a partition of $ [a,b]   $ and $ m = |P| - 1 $. As
\begin{equation*}
    \sum_{k=1}^{n_{j}} |f_{j}(x_{k})-f_{j}(x_{k-1}) | \leq \sum_{k=1}^{m} |f_{j}(x_{k})-f_{j}(x_{k-1}) | ,
\end{equation*}
then 
\begin{equation*}
    V(f,[a,b]) =  \sum_{j=1}^{n_{0}}[ V(f_{j},[a,b]) -\epsilon ] <   \sum_{j=1}^{n_{0}}  \sum_{k=1}^{n_{j}} |f_{j}(x_{k})-f_{j}(x_{k-1}) | \leq  \sum_{j=1}^{n_{0}}  \sum_{k=1}^{m} |f_{j}(x_{k})-f_{j}(x_{k-1}) | .
\end{equation*}

On the other hand, 

\thispagestyle{empty}
\newpage%
	
%%%%%%%%%%%%%%%%%%%%%%%%%%%%%%%%%%%%%%%%%%%%%%%%%%%%%%%%%%%%%%%%%%%%%%%%%%%%%%%%%%%%%%%%%%%%%%%%%%%%%%%%%%%%%%%%%%%%%%%%%%%%%%%%%%%%%%%%%%%%%%%%%%%%%%%%%%%%%%%%%%%%%%%%%%%%%%%%%%%%%%%%%%%%%%%%%%%%%%%%%%%%%%%%%%%%%%%%%%%%%%%%%%%%%%%%%%%%%%%%%%%%%%%%%%%%%%%%%%%%%%%%%%%%%%


\begin{align*}
  \sum_{j=1}^{n_{0}}  \sum_{k=1}^{m} |f_{j}(x_{k})-f_{j}(x_{k-1}) | & = \sum_{k=1}^{m} \sum_{j=1}^{n_{0}} |f_{j}(x_{k})-f_{j}(x_{k-1}) | \\
    & \leq \sum_{k=1}^{m} \sum_{j=1}^{\infty} |f_{j}(x_{k})-f_{j}(x_{k-1}) | \\
    & \leq V(f,P) \leq V(f, [a,b]) .
\end{align*}
So, $ V(f, [a,b]) < V(f, [a,b])   $ which is not possible.

Thus, 	$ \sum_{j=1}^{\infty} V(f_{j},[a,b])\leq V(f,[a,b])$.



	
	Now we will prove $  V(f,[a,b]) \leq  \sum_{j=1}^{\infty} V(f_{j},[a,b])\ $.\\

 \thispagestyle{empty}
\newpage%
	
%%%%%%%%%%%%%%%%%%%%%%%%%%%%%%%%%%%%%%%%%%%%%%%%%%%%%%%%%%%%%%%%%%%%%%%%%%%%%%%%%%%%%%%%%%%%%%%%%%%%%%%%%%%%%%%%%%%%%%%%%%%%%%%%%%%%%%%%%%%%%%%%%%%%%%%%%%%%%%%%%%%%%%%%%%%%%%%%%%%%%%%%%%%%%%%%%%%%%%%%%%%%%%%%%%%%%%%%%%%%%%%%%%%%%%%%%%%%%%%%%%%%%%%%%%%%%%%%%%%%%%%%%%%%%%
	Let $P=\{a=t_{0}<\cdots <t_{n}=b\}$ be a partition of $[a,b]$.
	Then 
	\begin{align*}
		\sum_{j=1}^{n}\|f(t_{j})-f(t_{j-1})\|_{1} & =\sum_{j=1}^{n}\left(\sum_{i=1}^{\infty}|f_{i}(t_{j})-f_{i}(t_{j-1})|\right)\\
		& =\sum_{i=1}^{\infty}\left(\sum_{j=1}^{n}|f_{i}(t_{j})-f_{i}(t_{j-1})|\right) \\
		& \leq \sum_{i=1}^{\infty}V(f_{i},[a,b]).
	\end{align*}
	
	
	
	
	Thus, $f$ is bounded variation and 
	
	\[V(f,[a,b])\leq \sum_{i=1}^{\infty}V(f_{i},[a,b]).\]
	
\thispagestyle{empty}
	\newpage
%%%%%%%%%%%%%%%%%%%%%%%%%%%%%%%%%%%%%%%%%%%%%%%%%%%%%%%%%%%%%%%%%%%%%%%%%%%%%%%%%%%%%%%%%%%%%%%%%%%%%%%%%%%%%%%%%%%%%%%%%%%%%%%%%%%%%%%%%%%%%%%%%%%%%%%%%%%%%%%%%%%%%%%%%%%%%%%%%%%%%%%%%%%%%%%%%%%%%%%%%%%%%%%%%%%%%%%%%%%%%%%%%%%%%%%%%%%%%%%%%%%%%%%%%%%%%%%%%%%%%%%%%%%%%%	

\thispagestyle{empty}
	\newpage

 
    \thispagestyle{empty}
	\newpage

 
    \begin{proof}[My proof]
        A proof.
    \end{proof}
    \begin{example}[My example]
        An example.
    \end{example}

   % \begin{block}{My block}
    %    A block.
    %\end{block}    



\thispagestyle{empty}
	\newpage
    \begin{frame}{El espacio $ l^{1} $}
    Propiedades:
    \begin{enumerate}
        \item Es un espacio de Schur.
        \item Todos los intervalos ordenados cerrados son compactos
        %\item Dependiendo su orden es normal o no. 
    \end{enumerate}
\end{frame}

\begin{frame}{Espacio $ C^{1} ([a,b], X) $}
Sea $X$ un ENO  donde $ X_{+} $ es normal y $ E^{1}  = C^{1} ( [a,b],X) $. Se define $E_{+}^{1} = \{ f \in E: f(t) \geq 0, f'(t) \geq 0 \ \forall t \in [a,b]  \} $ y $ || f ||_{\infty} = \sup \{ || f(t) ||_{X} : t \in [a,b] \} $.\\
      Propiedades:
    \begin{enumerate}
        \item Si $X_{+}$ es normal, entonces es normal.
        \item Si en $X$ todos los intervalos ordenados cerrados son compactos, entonces  en  todos los intervalos ordenados cerrados son compactos.
        \item Dependiendo su orden es normal o no. 
    \end{enumerate}
\end{frame}

\begin{frame}{Teoremas de Punto Fijo}
    \begin{block}{Teorema 1}
        Sean $X$ un ENO con la propiedad de Schur con $ X_{+} $ y $u_0, v_0 \in X$ con $u_0 < v_0$. \\
	Si  $A:[u_0,v_0] \rightarrow [u_0,v_0]$ subcontinuo creciente y existe $B$ débilmente compacto tal que $A([u_0, v_0]) \subseteq B$, entonces $A$ tiene un punto fijo minimal y maximal en $[u_0,v_0]$.
    \end{block}
    \begin{block}{Teorema 2}
        	Sean $X$ un ENO con la propiedad de que todo intervalo ordenado $ [a,b] $ es compacto y $u_0, v_0 \in X$ con $u_0 < v_0$. \\
	Si  $A:[u_0,v_0] \rightarrow [u_0,v_0]$  creciente, entonces $A$ tiene un punto fijo minimal y maximal en $[u_0,v_0]$.
    \end{block}
\end{frame}



\end{document}